% from ega4-2-fr.tex
\oldpage[IV-2]{192}

\begin{corollary}[7.2.7]
\label{IV.2.s7.2.7-fr}
Soit $A$ un anneau local noethérien tel qu'il existe un $A$-module $M$ de
type fini qui soit un $A$-module de Cohen--Macaulay et pour lequel
$\operatorname{Supp}(M)=\operatorname{Spec}(A)$; alors $A$ est strictement
formellement caténaire.
\end{corollary}

\begin{proof}
En effet, $A$ est formellement caténaire (7.1.5 et 7.1.9), et d'autre part
les fibres du morphisme canonique
$\operatorname{Spec}(\widehat A)\to\operatorname{Spec}(A)$ sont des
préschémas de Cohen--Macaulay (6.3.8), donc \emph{a fortiori} vérifient
$(\mathrm S_1)$.
\end{proof}

\begin{corollary}[7.2.8]
\label{IV.2.s7.2.8-fr}
Si $A$ est un anneau local noethérien strictement formellement caténaire,
tout anneau local $B$ qui est une $A$-algèbre essentiellement de type fini
est strictement formellement caténaire.
\end{corollary}

\begin{proof}
En effet, $B$ est formellement caténaire (7.1.11) et il résulte de
(7.4.4)\footnote{Le cor. (7.2.8) n'est pas utilisé pour démontrer (7.4.4).}
que les fibres de
$\operatorname{Spec}(\widehat B)\to\operatorname{Spec}(B)$ vérifient
$(\mathrm S_1)$, d'où la conclusion.
\end{proof}

\begin{corollary}[7.2.9]
\EGAReaderTarget{ega4.c04.cor.7-2-9}{IV.2.s7.2.9-fr}
Tout anneau local noethérien de dimension 1 est strictement formellement
caténaire.
\end{corollary}

\begin{proof}
En effet, tout anneau quotient intègre $B$ d'un tel anneau $A$ est de
dimension 0 ou 1, donc son complété est strictement équidimensionnel
(7.2.1.1).
\end{proof}

\begin{remarks}[7.2.10]
\label{IV.2.s7.2.10-fr}
\begin{enumerate}
\item[(i)] Il résulte de (7.2.7) et (7.2.8) que tout anneau quotient d'un
anneau local de Cohen--Macaulay est strictement formellement caténaire.

Un anneau local noethérien de dimension 2 vérifiant $(\mathrm S_2)$ est
strictement formellement caténaire, puisque c'est un anneau de
Cohen--Macaulay. Rappelons par contre qu'il y a des anneaux locaux
noethériens intègres de dimension 2 qui ne sont pas universellement
caténaires, ni \emph{a fortiori} strictement formellement caténaires
(5.6.11).
\item[(ii)] On ignore si un anneau formellement caténaire (7.1.9) est
strictement formellement caténaire; cela tient à ce qu'on ne sait pas si
pour un anneau local noethérien intègre $B$, $\widehat B$ est sans idéaux
premiers associés immergés (6.4.3). Nous ignorons également si, dans les
conditions équivalentes c), d), e) de (7.2.5), on peut remplacer
l'hypothèse que $A$ est universellement caténaire par l'hypothèse que $A$
est caténaire; on peut montrer qu'il en est ainsi lorsque $A$ est hensélien
(18.9.6).
\end{enumerate}
\end{remarks}

\subsection{Fibres formelles des anneaux locaux noethériens}
\label{subsection:IV.7.3-fr}

\begin{env}[7.3.1]
\label{IV.2.s7.3.1-fr}
Nous allons considérer dans ce numéro et les deux suivants des propriétés
$\boldsymbol P(Z,k)$ de la forme suivante :
\begin{quote}
« $Z$ est un préschéma localement noethérien sur un corps $k$,

et pour tout $z\in Z$, on a $\boldsymbol Q(\mathscr O_z,k)$ »
\end{quote}
où $\boldsymbol Q(A,k)$ est une propriété d'un anneau local noethérien $A$
qui est une $k$-algèbre; on suppose que si $k'$ est un corps isomorphe à
$k$, et si $A'$ est un anneau local noethérien qui est une $k'$-algèbre
di-isomorphe à $A$, les propriétés $\boldsymbol Q(A,k)$ et
$\boldsymbol Q(A',k')$ sont équivalentes.
\end{env}

